\section{Experimental Pipeline}
% ==========================================================
% Survey Construction and Generation Procedure (LaTeX)
% ==========================================================
\renewcommand{\appendixname}{}
Appendix A specifies the executable study procedure and supersedes any conflicting
procedural description in the main text unless explicitly noted.
\subsection*{Step 0 --- Discovery and shortlisting of the 25 cases}
The 25-case set was not assembled randomly. It was built through an iterative discovery and screening process
designed to produce a broad, information-rich evaluation set for structured analysis.

\subsubsection*{0.1 Initial discovery pool}
Candidate cases were initially identified from well-known historical failures and controversies across multiple
domains, including:
\begin{itemize}
  \item finance
  \item aerospace
  \item energy
  \item transportation
  \item software / technology
  \item corporate governance
  \item public-sector and institutional systems
\end{itemize}
The goal at this stage was breadth, not immediate inclusion.

\subsubsection*{0.2 Shortlisting criteria}
Candidates were screened against four requirements:
\begin{enumerate}
  \item \textbf{Major failure or major breakdown.} The case had to involve substantial financial loss, loss of life,
        mission failure, systemic disruption, or major governance breakdown.
  \item \textbf{Publicly describable system setup.} There had to be enough public information to describe the system,
        organizational structure, and task intent without relying on the final investigation findings.
  \item \textbf{Independent assessor material available.} The case had to have at least one credible external assessment source,
        such as:
        \begin{itemize}
          \item court findings
          \item congressional / parliamentary reports
          \item regulator findings
          \item independent investigation boards
          \item official commissions
        \end{itemize}
  \item \textbf{Usefulness for comparative analysis.} The case had to be suitable for comparing:
        \begin{itemize}
          \item baseline LLM analysis
          \item protocol-guided LLM analysis
          \item independent assessor reconstruction
        \end{itemize}
\end{enumerate}

\subsubsection*{0.3 Solution-space coverage}
During shortlisting, cases were chosen to span a wide solution space rather than cluster around a single failure type.
The shortlist was intentionally balanced across:
\begin{itemize}
  \item mechanistic engineering failures
  \item governance-heavy socio-technical failures
  \item financial and institutional failures
  \item regulatory and oversight failures
  \item software / information-system failures
\end{itemize}
This was done to ensure that the final survey could test whether structured reasoning helps more in some classes of cases
than others.

\subsubsection*{0.4 Exclusion criteria}
Cases were excluded if they failed one or more of the following:
\begin{itemize}
  \item no credible independent assessor source could be identified
  \item public documentation of system setup was too thin
  \item the case was too redundant with another selected case
  \item the case was still too unresolved or factually unstable for use as a frozen survey item
\end{itemize}

\subsubsection*{0.5 Finalization}
After iterative screening, 25 cases were selected and assigned stable IDs C01 through C25. The finalized list was then
frozen in \texttt{case\_list\_v1.txt}. This file became the root index for all subsequent survey construction steps.

\subsection*{Step 1 --- Construct the 25-case list}
The survey case set was constructed using the following criteria.

\paragraph{Selection requirements.}
Each case had to satisfy:
\begin{enumerate}
  \item Major system failure or controversy
  \item Availability of public documentation describing the system setup
  \item Availability of an independent investigation or assessment report
  \item Coverage across multiple socio-technical domains
\end{enumerate}

\paragraph{Domains intentionally covered.}
\begin{itemize}
  \item finance
  \item aerospace
  \item energy
  \item transportation
  \item technology / software
  \item governance / institutions
\end{itemize}
The objective was to create a heterogeneous solution space rather than a domain-specific dataset.

\subsection*{Step 2 --- Assign case identifiers}
Each case was assigned a stable identifier: \texttt{C01 \dots C25}. The case list was frozen in
\texttt{case\_list\_v1.txt}.

\paragraph{Example entries.}
\begin{itemize}
  \item \texttt{C01 --- London Whale (JPMorgan CIO)}
  \item \texttt{C02 --- Deepwater Horizon (Macondo Well)}
  \item \texttt{C03 --- Fukushima Daiichi Nuclear Plant}
  \item \texttt{\dots}
  \item \texttt{C25 --- Lehman Brothers Collapse}
\end{itemize}
This file is included in the freeze manifest.

\subsection*{Step 3 --- Create neutral case descriptions}
Neutral case descriptions were created by:
\begin{enumerate}
  \item taking the 25-case list
  \item writing $\leq 150$-word system descriptions
  \item removing failure conclusions and investigation findings
  \item freezing them in \texttt{neutral\_case\_descriptions\_v1.txt}
  \item converting them in the notebook into per-case inputs:
        \texttt{Cxx\_case\_input.txt}
\end{enumerate}
These descriptions describe system structure and task intent, not the failure outcome.

\paragraph{Historical context constraint.}
Historical context may be included in a neutral description only when necessary to define the operational setting,
system boundary, or decision environment of the case. Such context must be stated as descriptive background facts and
must not introduce causal interpretations, evaluative judgments, or conclusions from investigation reports. Any historical
material included must be provided identically to all model-generation conditions.

\subsection*{Step 4 --- Freeze survey inputs}
The following artifacts were frozen with SHA-256 hashes:
\begin{itemize}
  \item \texttt{case\_list\_v1.txt}
  \item \texttt{neutral\_case\_descriptions\_v1.txt}
  \item \texttt{prompt\_baseline\_v1.txt}
  \item \texttt{prompt\_protocol\_v1.txt}
  \item \texttt{prompt\_independent\_v1.txt}
  \item \texttt{protocol\_specification\_v1.txt}
  \item \texttt{model\_config\_v1.json}
  \item \texttt{case\_id\_mapping\_v1.csv}
\end{itemize}
These are recorded in \texttt{case\_freeze\_manifest\_v1.txt}.

\subsection*{Step 5 --- Notebook input construction}
The survey notebook performs the following preprocessing:
\[
\texttt{neutral\_case\_descriptions\_v1.txt}
\;\;\rightarrow\;\;
\text{parse by case ID}
\;\;\rightarrow\;\;
\text{write case inputs}
\;\;\rightarrow\;\;
\texttt{C01\_case\_input.txt, \dots, C25\_case\_input.txt}.
\]
These files are the primary input artifacts for model generation.

\subsection*{Step 6 --- Generation conditions}
Three analysis conditions are defined.

\paragraph{Baseline (B).}
\textbf{Inputs:} baseline prompt + case input.\\
\textbf{Output:} \texttt{Cxx\_B\_raw.txt}.

\paragraph{Protocol (P).}
\textbf{Inputs:} protocol prompt + \texttt{protocol\_specification\_v1.txt} + case input.\\
\textbf{Output:} \texttt{Cxx\_P\_raw.txt}.

\paragraph{Independent (I).}
\textbf{Inputs:} independent prompt + case input + assessor report extracts.\\
\textbf{Outputs:} \texttt{Cxx\_I\_refA\_raw.txt}, \texttt{Cxx\_I\_refB\_raw.txt}.\\
Extract files are automatically constructed using an LLM via near-verbatim, source-bounded extraction from official investigation reports, with human review limited to the E/T boundary (Section 7.4.1).
\paragraph{Condition discipline separation (normative).}
Only the Protocol condition (P) is required to follow the FRFP-compatible executable
instruction surface. Baseline (B) represents unconstrained model analysis.
Independent (I) differs from (B) by information access (assessor extracts), not by
protocol discipline, unless a separate hybrid condition is explicitly defined.

\subsection*{Step 7 --- Selection and acquisition of Independent Assessor Reports and supporting materials 
The Independent analysis condition requires access to information derived from credible external investigations of each case.
A structured procedure is used to (i) identify authoritative assessment sources and (ii) acquire the corresponding source files
as frozen local snapshots.

\subsubsection*{7.1 Source hierarchy}
For each case, candidate assessor sources are evaluated according to a predefined credibility hierarchy:
\begin{enumerate}
  \item Court findings or judicial determinations
  \item Congressional or parliamentary investigation reports
  \item Regulatory authority reports
  \item Independent investigation boards or commissions
\end{enumerate}

\subsubsection*{7.2 Primary, alternate, and supplemental assessor sources}
For each case, assessor sources are recorded in a canonical JSON log:
\begin{itemize}
  \item \textbf{Primary A} --- the most authoritative available assessment according to the hierarchy above.
  \item \textbf{Primary B} --- an alternate independent assessment source for sensitivity analysis (when available).
  \item \textbf{Supplemental} --- an additional technical/operational source used only when it contributes material not already covered by the primaries.
\end{itemize}

\subsubsection*{7.3 Supporting materials (discovery and selection rules)}
Supporting materials are \emph{inputs to} an alternate evidence-augmented analysis and are treated as upstream evidence that assessor reports analyze.
Accordingly:
\begin{enumerate}
  \item \textbf{Rule (main-report exclusion):} The \emph{main assessor report text} is \emph{not} included as supporting material.
        Supporting materials must be distinct upstream artifacts (e.g., testimony, exhibits, technical appendices) that the main report relies on.
  \item \textbf{Rule (official channel preference):} Supporting materials are selected only if present or acquirable via an official inquiry / investigation channel
        (e.g., inquiry document repositories, official hearing transcript archives, regulator evidence dockets).
  \item \textbf{Rule (use-driven selection):} Materials are chosen to support an alternate AI analysis grounded in primary evidence (not the report’s conclusions).
        Preference is given to records that expose system architecture, operational process, decision structure, and causal evidence.
  \item \textbf{Allowed material types:} Each supporting item is labeled with \texttt{material\_type} from:
        \{\texttt{transcript}, \texttt{submission}, \texttt{evidence}, \texttt{hearing}, \texttt{other}\}.
\end{enumerate}

\subsubsection*{7.4 Acquisition, storage, and manifesting (implemented)}
This step is implemented as a deterministic acquisition pipeline driven by \\ \texttt{independent\_assessor\_source\_log\_v1.json}
and \texttt{extraction\_config.json}:
\begin{enumerate}
  \item Download snapshots of \texttt{PrimaryA}, \texttt{PrimaryB} (if present), \texttt{Supplemental} (if present) into:
        \texttt{independent\_sources/sources/Cxx/\{primaryA,primaryB,supplemental\}/}.
  \item Download supporting materials into:
        \texttt{independent\_sources/sources/Cxx/supporting\_materials/}.
  \item Maintain a single \texttt{source\_manifest.csv} that records, for every downloaded item:
        \texttt{case\_id, role, title, url, hierarchy, supplemental\_type, material\_type, retrieved\_utc, local\_path, sha256}.
  \item A manifest synchronization pass ensures downloaded files and the manifest remain consistent (including filename normalization and missing-row repair).
\end{enumerate}

\subsection*{Step 8 --- LLM generation of analysis artifacts (B, P, I)}
After the dataset and input artifacts were frozen, a controlled generation procedure was used to produce analytical outputs
from the language model under three different reasoning conditions. The objective was to produce comparable analytical artifacts
that differ only in the reasoning constraints and information provided to the model.

\subsubsection*{8.1 Model configuration}
All generations were performed using a fixed model configuration defined in \texttt{model\_config\_v1.json}. The configuration specifies:
\begin{itemize}
  \item model provider
  \item model version
  \item temperature
  \item output token limits
\end{itemize}
The configuration file was frozen alongside the dataset to ensure reproducibility.
In the implemented pipeline, generation conditions differ by which locally downloaded materials are summarized and injected.


\subsubsection*{8.2 Input construction}
For each case identifier \texttt{Cxx}, the notebook constructs a case input file \texttt{Cxx\_case\_input.txt}.
This file contains the neutral case description corresponding to that case. The case input files serve as the base input
across all three reasoning conditions.

\subsubsection*{8.3 Generation conditions}
Inputs: \texttt{prompt\_baseline\_v1.txt} + \texttt{Cxx\_case\_input.txt} + summarized supporting materials (if available).\\
Output: \texttt{Cxx\_B\_raw.txt}.\\
If no supporting materials exist for a case, Baseline generation for that case is skipped.


\subsubsection*{8.4 Protocol condition (P)}
Inputs: \texttt{prompt\_protocol\_v1.txt} + \texttt{protocol\_specification\_v1.txt} + \texttt{Cxx\_case\_input.txt}
+ summarized supporting materials (if available).\\
Output: \texttt{Cxx\_P\_raw.txt}.\\
If no supporting materials exist for a case, Protocol generation for that case is skipped.

The protocol specification defines procedural constraints such as intent locking, artifact discipline, branch tracking,
explicit assumptions, and separation of analysis and endorsement.

\paragraph{Protocol handoff is specified but not executed (normative).}
In Condition P, the handoff/verification component required by the FRFP-compatible
instruction surface is produced only as an explicit artifact (a proposed review and
verification plan consistent with the $E \rightarrow T$ boundary).
No post-generation human endorsement, correction, or modification of the model output
is performed within the experiment.

\subsubsection*{8.5 Independent condition (I)}
Inputs: \texttt{prompt\_independent\_v1.txt} + \texttt{Cxx\_case\_input.txt} + summarized assessor sources:
\texttt{PrimaryA (+ Supplemental)} for \texttt{I\_refA}, and \texttt{PrimaryB (+ Supplemental)} for \texttt{I\_refB}.\\
Outputs: \texttt{Cxx\_I\_refA\_raw.txt}, \texttt{Cxx\_I\_refB\_raw.txt}.

\subsubsection*{8.6 Execution procedure}
The generation notebook iterates over case identifiers and executes:
\[
\text{Baseline (B)} \rightarrow \text{Protocol (P)} \rightarrow \text{Independent (I\_refA)} \rightarrow \text{Independent sensitivity (I\_refB)}.
\]
Each generation produces a raw text artifact stored in the experiment output directory.
\subsubsection*{Style isolation rule.}

To preserve separation between agent configurations, stylistic conventions
introduced by the protocol condition must not appear in Baseline or
Independent outputs. In particular, labeled reasoning artifacts used by the
protocol workflow (e.g., \texttt{FACT:}, \texttt{ASSUMPTION:}, \texttt{MECHANISM:},
\texttt{CONTROL:}, \texttt{RISK:}, \texttt{UNCERTAINTY:}) are prohibited in
Baseline and Independent generations.

During contamination checks prior to normalization, outputs are inspected for
cross-condition stylistic leakage. If protocol-specific artifact labels or
reasoning structures appear in Baseline or Independent outputs, the affected
case is regenerated under the isolation protocol and the discarded artifacts
are logged.
\subsubsection*{8.7 Output logging and hashing}
For every generated output:
\begin{enumerate}
  \item The raw output text is saved to disk.
  \item A SHA-256 hash of the file is computed.
  \item A generation log entry is recorded containing:
        \begin{itemize}
          \item case identifier
          \item generation condition
          \item model configuration
          \item timestamp
          \item output file hash
        \end{itemize}
\end{enumerate}
These records ensure that all outputs can be verified against the frozen experimental configuration.

\subsubsection*{8.8 Resulting artifact set}
For each case, the generation step produces:
\begin{itemize}
  \item \texttt{Cxx\_B\_raw.txt}
  \item \texttt{Cxx\_P\_raw.txt}
  \item \texttt{Cxx\_I\_refA\_raw.txt}
  \item \texttt{Cxx\_I\_refB\_raw.txt}
\end{itemize}
In the full experiment with 25 cases, this yields up to $25 \times 4 = 100$ generated artifacts. These raw outputs form the basis
for subsequent normalization and human evaluation stages.

\subsection*{Step 9 --- Output normalization and blinding}
Raw outputs differ in length, style, and structure across reasoning conditions. To ensure fair comparison during human evaluation,
all outputs were normalized into a consistent presentation format and blinded with respect to their generation condition.
Normalization does not alter substantive analytical content.

\subsubsection*{9.1 Objective of normalization}
Normalization serves three purposes:
\begin{enumerate}
  \item \textbf{Structural comparability:} ensure outputs can be evaluated on the same criteria.
  \item \textbf{Removal of formatting differences:} eliminate stylistic cues that might reveal the generation condition.
  \item \textbf{Preparation for human rating:} convert outputs into standardized one-page analytical summaries suitable for pairwise evaluation.
\end{enumerate}

\subsubsection*{9.2 Standard output structure}
Each raw output was converted into a standardized analytical template consisting of:
\begin{enumerate}
  \item System Objective
  \item Representation and Models
  \item Measurement and Information Flows
  \item Control Mechanisms
  \item Decision Interfaces
  \item Governance and Oversight
  \item Possible Failure Mechanisms
  \item Candidate Interventions or Safeguards
\end{enumerate}

\subsubsection*{9.3 Normalization process}
The normalization procedure consisted of:
\begin{enumerate}
  \item \textbf{Section extraction:} parse raw output to identify elements corresponding to the standardized sections.
  \item \textbf{Reformatting:} reorganize extracted content into the standardized template structure.
  \item \textbf{Length harmonization:} trim extremely long responses while preserving key analytical points.
  \item \textbf{Removal of condition-specific cues:} remove references that might reveal the generation condition (e.g., ``protocol'',
        ``baseline'', investigation reports, or internal reasoning procedures).
\end{enumerate}
\paragraph{Modality preservation (normative).}
Normalization must preserve epistemic modality and endorsement stance (e.g., expressions of certainty, hedging, and
evidence-vs-assumption separation) as substantive content. Only condition-revealing labels (e.g., ``baseline'',
``protocol'', or explicit references to assessor sources) and superficial formatting cues may be removed.
\subsubsection*{9.4 Blinding procedure}
After normalization, each document is assigned an anonymized identifier (e.g., \texttt{D0001}, \texttt{D0002}, \dots). The mapping between
document identifiers and generation conditions is stored in a separate key file not visible to raters. This key file is retained by the
experiment administrator and is not included in rater packets.

\subsubsection*{9.5 Pair construction for rating}
Raters evaluate documents via pairwise comparisons. For each case, the following comparisons are constructed:
\begin{itemize}
  \item Baseline vs Protocol
  \item Protocol vs Independent
  \item Baseline vs Independent
\end{itemize}
Pairs are randomized across raters. Each comparison consists of two normalized one-page documents labeled \textbf{Document A} and \textbf{Document B},
with ordering randomized to prevent position bias.
\subsubsection*{9.X Per-rater Independent reference selection}
Pair construction is performed independently per rater. For each case requiring an Independent comparison,
the rater-specific pairing logic may select \texttt{I\_refA} or \texttt{I\_refB}. A fixed random seed is used to make the
per-rater selection deterministic and reproducible. The combined mapping is stored in \texttt{pairs\_key\_all\_raters.csv}.

\subsubsection*{9.Y Main-analysis restriction }
For the main analysis dataset, comparisons involving the Independent condition are filtered to include only
pairs where the Independent document is \texttt{I\_refA}. Pairs containing \texttt{I\_refB} are reserved for sensitivity analysis.
\subsubsection*{9.6 Blinded rater packets}
Each rater receives a packet containing:
\begin{itemize}
  \item randomized document pairs
  \item evaluation instructions
  \item scoring rubric
\end{itemize}
Packets contain no information about generation condition, case identifiers, protocol usage, or investigation sources.

\subsubsection*{9.7 Preservation of raw outputs}
All raw model outputs remain archived without modification:
\texttt{Cxx\_B\_raw.txt}, \texttt{Cxx\_P\_raw.txt}, \texttt{Cxx\_I\_refA\_raw.txt}, \texttt{Cxx\_I\_refB\_raw.txt}.
Normalization produces derived evaluation documents, while raw artifacts remain the canonical record of model output.

\subsection*{Step 11 --- Statistical analysis and comparison methodology}
The objective is to determine whether outputs generated using the structured reasoning protocol demonstrate improved analytical quality
relative to baseline outputs. Because documents are evaluated using pairwise comparisons, analysis focuses on relative preference frequencies.

\subsubsection*{11.1 Pairwise comparison structure}
For each case, three comparisons are constructed: Baseline vs Protocol, Protocol vs Independent, Baseline vs Independent.
Each comparison is evaluated across four rubric dimensions:
\begin{itemize}
  \item \textbf{R1} --- Analytical clarity
  \item \textbf{R2} --- Structural reasoning quality
  \item \textbf{R3} --- Diagnostic usefulness
  \item \textbf{R4} --- Overall analytical usefulness
\end{itemize}
For each dimension, raters select: \emph{A closer}, \emph{B closer}, or \emph{No clear difference}.

\subsubsection*{11.2 Data representation}
Each rating is represented as a record containing:
\[
\texttt{case\_id, document\_A\_id, document\_B\_id, generation\_condition\_A, generation\_condition\_B, rater\_id, dimension, response},
\]
where \texttt{response} is one of \texttt{A}, \texttt{B}, or \texttt{Tie}.

\subsubsection*{11.3 Preference aggregation}
For each comparison type and dimension, counts are computed:
$N(A>B)$, $N(B>A)$, and $N(\mathrm{Tie})$.
Tie responses are excluded from directional preference calculations but retained in reporting.
The primary statistic is the preference rate:
\[
\mathrm{PreferenceRate}(A>B) \;=\; \frac{N(A>B)}{N(A>B) + N(B>A)}.
\]

\subsubsection*{11.4 Hypothesis testing}
The principal hypothesis tested is:
\[
H_0:\ \text{Protocol outputs are no more likely to be preferred than Baseline outputs.}
\]
This is evaluated using a binomial test (or sign test) on directional preference counts. For each dimension:
\[
p\text{-value} \;=\; \mathrm{BinomialTest}\!\left(N(P>B),\ N(P>B)+N(B>P),\ p=0.5\right).
\]

\subsubsection*{11.5 Effect size estimation}
Effect size is reported using a preference difference:
\[
\mathrm{EffectSize} \;=\; \mathrm{PreferenceRate}(P>B) \;-\; \mathrm{PreferenceRate}(B>P).
\]
Values greater than zero indicate preference for the protocol condition.

\subsubsection*{11.6 Cross-case aggregation}
Results are analyzed at two levels:
\begin{itemize}
  \item \textbf{Case-level analysis:} preference statistics computed separately for each case.
  \item \textbf{Aggregate analysis:} preference counts pooled across all cases to estimate overall performance differences.
\end{itemize}

\subsubsection*{11.7 Sensitivity analysis}
To evaluate robustness of the Independent condition, comparisons are repeated using the alternate assessor source:
\texttt{Independent\_refA} vs \texttt{Independent\_refB}. Differences provide an estimate of sensitivity to assessor source selection.

\subsubsection*{11.8 Inter-rater consistency}
Inter-rater agreement is evaluated using measures such as Fleiss' kappa or pairwise agreement rates across raters.

\subsubsection*{Lexical contamination check.}
In addition to manual review, a simple lexical contamination check is performed
prior to normalization and packet construction. For each case, the raw outputs
from Baseline, Protocol, and Independent conditions are compared for leakage of
condition-specific marker terms.

The check includes:
\begin{itemize}
  \item protocol-unique artifact labels (e.g., \texttt{FACT:}, \texttt{ASSUMPTION:},
        \texttt{MECHANISM:}, \texttt{CONTROL:}, \texttt{RISK:}, \texttt{UNCERTAINTY:});
  \item protocol-unique structural terms (e.g., branch labels, handoff language,
        explicit verification-path language);
  \item forbidden theoretical vocabulary in Baseline or Independent outputs.
\end{itemize}

For each raw output, the contamination script records:
\begin{itemize}
  \item presence/absence of protected marker terms;
  \item overlap counts for condition-specific lexicons;
  \item flagged cases requiring manual review.
\end{itemize}

A case is marked for regeneration if Baseline or Independent outputs contain
protocol-specific artifact labels or other condition-specific lexical markers
that would compromise workflow isolation. All flagged cases, regeneration
decisions, and discarded artifact hashes are recorded in the experiment log.

\subsubsection*{11.9 Interpretation}
The primary outcome of interest is whether the Protocol condition demonstrates:
\begin{itemize}
  \item higher preference rates than the Baseline condition,
  \item consistent improvement across evaluation dimensions,
  \item robustness across cases and raters.
\end{itemize}
The Independent condition serves as a reference representing analyses grounded in formal investigation reports.